\section[Criterios estructurales y carta metrológica de \(G\)]{Criterios de
selección de la coordenada gravitatoria: holonomía enriquecida, respuesta
espectral y carta metrológica}
\label{sec:selector-gravitatorio-rev11}
\index{constante gravitatoria!transductor y selector}
\index{gravedad!coordenada metrológica}
\index{holonomía!selector gravitatorio}

La composición APP--TRIT--TPK, el estado enriquecido y la rama de acción
construyen primero la coordenada gravitatoria. La constante gravitatoria se
organiza en tres niveles tipados: la recta dimensional que alberga \(G\), el
transductor que convierte una coordenada adimensional en una sección de esa
recta y el selector estructural de dicha coordenada. La acción
Palatini--Cartan--Holst y su Hamiltoniano restringido realizan después la
salida ya generada. La carta decimal y la incertidumbre experimental se
presentan al final como representación y contraste.

\subsection{Recta dimensional y transductor interno}

En la retícula dimensional de la Mecánica Dimensional,
\[
 \mathbf G=-\mathbf S+5\mathbf C+2\mathbf T,
 \qquad [G]=L^3M^{-1}T^{-2}.
\]
Sean
\[
 c_{\rm int}\in\mathcal L_C^+,
 \qquad t_0\in\mathcal L_T^+,
 \qquad \ell_0=c_{\rm int}t_0,
 \qquad \hbar_{\rm int}\in\mathcal L_S^+.
\]

\HMTClaim{MD-R-G-TRANSDUCER}
\begin{definition}[Transductor gravitatorio]
\label{def:transductor-gravitatorio-rev11}
Para \(\theta_G>0\), se define
\begin{equation}
 \boxed{
 G_{\rm int}(\theta_G)
 =\theta_G\frac{c_{\rm int}^{5}t_0^2}{\hbar_{\rm int}}
 =\theta_G\frac{c_{\rm int}^{3}\ell_0^2}{\hbar_{\rm int}}
 \in\mathcal L_G^+.}
 \label{eq:transductor-gravitatorio-rev11}
\end{equation}
\end{definition}

La igualdad de las dos expresiones y el tipo de
\eqref{eq:transductor-gravitatorio-rev11} son exactos. Si
\[
 L_G^2:=\frac{\hbar_{\rm int}G_{\rm int}}{c_{\rm int}^3},
\]
entonces
\begin{equation}
 \theta_G=\left(\frac{L_G}{\ell_0}\right)^2.
 \label{eq:thetaG-razon-longitudes-rev11}
\end{equation}
La coordenada gravitatoria es, por tanto, una razón cuadrática de longitudes
o una respuesta de holonomía normalizada. La década
\(10^{-34}\mathcal U_S\) ya pertenece a \(\hbar_{\rm int}\): fija la sección
absoluta de acción y no constituye un factor relativo añadido a \(G\).

\subsection{Torsión observable y caracteres positivos}

El retorno observable de \(27\) pasos
\[
 \mathcal K_{27}:=F_{\rm obs}^{27}
\]
satisface
\[
 \mathcal K_{27}^{4}=F_{\rm obs}^{108}=I.
\]
En la coordenada \(\theta\in\mathbb Z/108\mathbb Z\), el orden es
exactamente cuatro, pues
\(\mathcal K_{27}(\theta)=\theta+27\),
\(\mathcal K_{27}^{2}(\theta)=\theta+54\neq\theta\) y
\(\mathcal K_{27}^{4}(\theta)=\theta\).

\HMTClaim{MD-R-G-POSITIVE-CHARACTER-NOGO}
\begin{theorem}[Trivialidad del carácter positivo observable]
\label{thm:caracter-positivo-G-rev11}
Todo homomorfismo multiplicativo
\[
 \chi:\langle\mathcal K_{27}\rangle
 \longrightarrow\mathbb R_{>0}^{\times}
\]
cumple \(\chi(\mathcal K_{27})=1\).
\end{theorem}

\begin{proof}
La multiplicatividad implica
\(\chi(\mathcal K_{27})^4=\chi(I)=1\).
El único real positivo cuya cuarta potencia es uno es \(1\).
\end{proof}

El mismo ciclo sí transporta los caracteres unitarios
\[
 \chi_j(\mathcal K_{27})=\exp(2\pi ij/4),
 \qquad j=0,1,2,3.
\]
El cociente observable admite fase, mientras que su torsión elimina una
escala positiva multiplicativa no trivial. Una continuación multiplicativa debe, por ello,
actuar sobre un levantamiento \(\widetilde{\mathcal K}_{27}\) cuya clase en
la abelianización del grupo de isotropía sea no torsional, y debe seleccionar
de manera natural tanto el carácter como su normalización.

\subsection{Memoria cronológica y respuesta espectral}
\label{sec:respuesta-espectral-G-rev11}

La cobertura cronológica dirigida realiza el retorno de \(27\) pasos como la
traslación
\[
 (G_{\rm or},S)\longmapsto(G_{\rm or},S)+(1,18),
\]
de orden semigrupal infinito. Esta vía conserva acumulación, pero su
respuesta espectral requiere interferencia entre historias.

Sea \(\Gamma=(V,E)\) un multigrafo dirigido finito con medida positiva,
pesos \(a_\ast(e)\geq0\) y cociclo \(c(e)\in\mathbb Z^2\). El operador
torcido
\[
 (\mathcal L_{\beta,\vartheta}f)(y)
 =\sum_{e:x\to y}
 e^{-\beta a_\ast(e)}e^{i\langle\vartheta,c(e)\rangle}f(x)
\]
induce el operador positivo
\(\mathcal H_{\beta,\vartheta}
=\mathcal L_{\beta,\vartheta}^{*}\mathcal L_{\beta,\vartheta}\).

\HMTClaim{MD-R-G-PHASE-INTERFERENCE}
\begin{proposition}[Criterio de interferencia para la respuesta de fase]
\label{prop:interferencia-selector-G-rev11}
En una órbita determinista con una única arista entrante por vértice,
\(\mathcal H_{\beta,\vartheta}\) es independiente de \(\vartheta\). Toda
respuesta espectral de fase distinta de cero requiere historias coherentes
múltiples o un diferencial magnético cuya holonomía sobreviva en los ciclos
enriquecidos.
\end{proposition}

\begin{proof}
En la órbita determinista, la única fase de cada término se multiplica por su
conjugada en
\(\mathcal L_{\beta,\vartheta}^{*}\mathcal L_{\beta,\vartheta}\)
y se cancela. La ausencia de términos cruzados elimina toda dependencia de
\(\vartheta\).
\end{proof}

Si una banda simple \(\lambda_0(\vartheta)\) posee un mínimo en el origen,
su Hessiano \(Q\) es semidefinido positivo. Con
\[
 v_{27}=(1,18),
 \qquad
 \mathcal R_{27}=v_{27}^{\mathsf T}Qv_{27},
\]
la condición \(\mathcal R_{27}\neq0\) constituye el control mínimo de
sensibilidad a la memoria acumulada. Una normalización natural
\(\mathfrak N\) permitiría definir
\[
 \Lambda_G=\ell_0\mathfrak N(Q,v_{27}),
 \qquad
 G_{\rm spec}
 =\frac{c_{\rm int}^{3}}{\hbar_{\rm int}}
  \bigl(\delta_\theta\Lambda_G\bigr)^2,
 \qquad
 \delta_\theta=\frac{A_{\rm deg}-C_{\rm deg}^{\ast}}{360}.
\]
La homogeneidad de esta plantilla es exacta. Su realización física queda
caracterizada por cuatro condiciones positivas: multigrafo enriquecido,
respuesta no nula, estabilidad por refinamiento e independencia de reescala,
junto con una normalización \(\mathfrak N\) fijada antes del contraste
metrológico.

\subsection{Funcional geométrico sobre la holonomía enriquecida}

La segunda continuación define directamente un funcional de longitud
\[
 \mathscr L:
 \operatorname{Loop}(\mathcal X_{\rm TPK}^{\rm enr})
 \longrightarrow\mathcal L_L^+\cup\{0\},
\]
invariante por cambio cíclico de punto base, conjugación, reetiquetado
admisible y subdivisión. Para un levantamiento enriquecido se escriben
\[
 \Lambda_{27}:=\mathscr L(\widetilde{\mathcal K}_{27}),
 \qquad
 \rho_{27}:=\frac{\Lambda_{27}}{c_{\rm int}t_0},
 \qquad
 \theta_G^{\rm hol}:=(\delta_\theta\rho_{27})^2.
\]
La transducción produce la plantilla homogénea
\begin{equation}
 \boxed{
 G_{\rm hol}
 =\frac{c_{\rm int}^{3}}{\hbar_{\rm int}}
  \bigl(\delta_\theta\Lambda_{27}\bigr)^2.}
 \label{eq:plantilla-holonomia-G-rev11}
\end{equation}
La ley de concatenación de este funcional sobre el ciclo torsional forma
parte de su definición y de su prueba. La selección física se completa
construyendo \(\widetilde{\mathcal K}_{27}\) y \(\mathscr L\), fijando una
normalización única y verificando la estabilidad de la salida. El valor
\(G_{\rm CODATA\,2022}\) interviene exclusivamente después, en el contraste
metrológico.

\begin{proposition}[Separación de los dos selectores gravitatorios]
\label{prop:separacion-selectores-G}
El selector circular
\(\theta_{108}^{\rm circ}=(54/\pi)^2\), construido por el cierre
acción--Compton--gravedad, y el selector holonómico
\(\theta_G^{\rm hol}=(\delta_\theta\rho_{27})^2\) son dos realizaciones
tipadas distintas del transductor \(G_{\rm int}\).  No se identifican entre
sí y la carta decimal \(\mathscr R_G\) no constituye una evaluación de
\(G_{\rm hol}\).
\end{proposition}
\begin{proof}
El primer selector depende únicamente del cierre circular de orden 108. El
segundo contiene, además, la longitud espectral
\(\Lambda_{27}=\mathscr L(\widetilde{\mathcal K}_{27})\).  Una igualdad
entre ambos requeriría calcular \(\Lambda_{27}\) y probar la coincidencia de
normalizaciones; ninguna operación del lector decimal realiza ese cálculo.
Por tanto, sus dominios y datos de entrada los distinguen antes de todo
contraste numérico.
\end{proof}

\subsection{Criterio de selección estructural}

\HMTClaim{MD-R-G-SELECTOR-CRITERION}
\begin{criterion}[Selector estructural de la constante gravitatoria]
\label{crit:selector-G-rev11}
Un selector HMT de \(G\) queda identificado cuando satisface conjuntamente:
\begin{enumerate}
 \item dominio enriquecido, composición y orientación declarados;
 \item descenso a una clase abeliana no torsional, o funcional geométrico
 invariante por punto base, conjugación, reetiquetado y subdivisión;
 \item respuesta positiva no trivial y estable bajo refinamiento;
 \item normalización única e independiente de la escala global del operador;
 \item selección interna de la coordenada y de su posición decimal;
 \item congelación del código y de las entradas antes de abrir el valor de
 contraste;
 \item realización final como sección de \(\mathcal L_G\) en una carta
 dimensional común.
\end{enumerate}
\end{criterion}

Las razones adimensionales
\[
 \frac{\mathcal G_{\rm tor}}{\mathcal I_{\rm tor}}=1,
 \qquad
 \frac{\mathcal G_{\rm av}}{\mathcal I_{\rm av}}
 =\frac{\pi_{\rm HMT}}{\varphi^2}
\]
describen lectores gravitatorio--inerciales distintos y permanecen
compatibles con una sola sección dimensional \(G\). La construcción de esa
sección utiliza además
\((\hbar_{\rm int},c_{\rm int},\ell_0)\) y el selector \(\theta_G\).

\subsection{Evaluación de la carta decimal gravitatoria y contraste metrológico posterior}

La representación decimal conservada en el corpus es
\begin{equation}
 \mathscr R_G([n,k,d]_3;t,e)
 :=10^{-n}\left(k+\frac{t}{10^3}+\frac{e}{10^d}\right).
 \label{eq:lector-decimal-G-rev11}
\end{equation}
Para la carta
\([n,k,d]_3=[11,6,5]_3\), \(t=674\) y \(e=30\), la evaluación racional es
\begin{equation}
 \boxed{
 \mathscr R_G
 =10^{-11}\left(6+\frac{674}{1000}+\frac{30}{10^5}\right)
 =6.67430\times10^{-11}.}
 \label{eq:evaluacion-decimal-G-rev11}
\end{equation}
La terna, los dos bloques y la posición decimal constituyen los datos
declarados de esta carta. Su evaluación es exacta; su selección estructural
queda regida por el \cref{crit:selector-G-rev11}.

Tras elegir la base SI de \(\mathcal L_G\), el contraste vigente es
\begin{equation}
 \boxed{
 G_{\rm CODATA\,2022}=6.67430(15)\times10^{-11}
 \ \mathrm{m^3\,kg^{-1}\,s^{-2}}.}
 \label{eq:G-CODATA-2022-rev11}
\end{equation}

\begin{center}
\begin{tabularx}{0.94\textwidth}{@{}L{0.27\textwidth}Y@{}}
\toprule
Campo & Declaración metrológica\\
\midrule
Valor central &
\(6.67430\times10^{-11}\ \mathrm{m^3\,kg^{-1}\,s^{-2}}\)\\
Incertidumbre estándar &
\(u(G)=0.00015\times10^{-11}
=1.5\times10^{-15}\ \mathrm{m^3\,kg^{-1}\,s^{-2}}\)\\
Incertidumbre relativa & \(u_r(G)\simeq2.2\times10^{-5}\)\\
Fuente y temporalidad &
Valores recomendados CODATA 2022, publicados en 2024
\cite{CODATA2022}\\
Función probatoria &
Contraste externo posterior; la unidad y la incertidumbre no intervienen en
la evaluación racional de \eqref{eq:evaluacion-decimal-G-rev11}.\\
\bottomrule
\end{tabularx}
\end{center}

\paragraph{Reproducibilidad del lector gravitatorio.}
Una enumeración independiente comprueba conjuntamente el tipo dimensional, el
lector decimal y las identidades de monodromía sobre el dominio finito
declarado.

\begin{proposition}[Cadena causal de la constante gravitatoria]
La acción covariante, el Hamiltoniano restringido y la ecuación
Einstein--Schrödinger gobiernan la dinámica gravitatoria. El tipo dimensional
y el transductor \eqref{eq:transductor-gravitatorio-rev11} son exactos; la
evaluación \eqref{eq:evaluacion-decimal-G-rev11} es exacta para su carta
declarada. El valor, la unidad y la incertidumbre de
\eqref{eq:G-CODATA-2022-rev11} constituyen contraste externo CODATA 2022. La
selección autónoma de \(\theta_G\) queda tipada por el levantamiento
enriquecido, la respuesta interferente y el funcional de holonomía sometidos
al \cref{crit:selector-G-rev11}.
\end{proposition}
